% MACRO I (ECO387C), Fall 2026 - mock midterm, exam only (answer space after
% every part, one question per page start). Solutions: mock-midterm-solutions.tex.
% Author: Philip Livdan. Four questions, each broken into lettered parts:
% Q1 stochastic endowment economy (PSet 6 with CRRA gamma=2 and asymmetric
% persistence), Q2 anticipated income rise in a two-type endowment economy
% (Lecture 2, textbook Sec 5.2.1 Arrow-Debreu prices), Q3 growth model by
% dynamic programming (textbook Examples 4.1, 4.5-4.8, Lecture 3), Q4
% consumption (textbook Sec 9.3.2-9.3.3, Lecture 4).
% Every number is checked in mock-midterm-checks.py (-checks.txt).
% Font gotchas: 12pt only, no math in \section*, \mathbf not \boldsymbol.
\documentclass[12pt]{article}
\usepackage{amsmath,amssymb}
\usepackage[margin=1in]{geometry}
\usepackage{needspace}
\DeclareMathSizes{12}{12}{8}{7}
\DeclareMathSizes{11}{11}{8}{7}
\DeclareMathSizes{10}{10}{8}{7}
\renewcommand{\ttdefault}{lmr}
\setlength{\parindent}{0pt}
\setlength{\parskip}{0.5\baselineskip}
\allowdisplaybreaks
\setlength{\emergencystretch}{3em}
\newcommand{\E}{E}
\newcommand{\gap}[1]{\par\vspace{#1}}
\newcommand{\qp}[2]{\par\medskip\textbf{(#1)} (#2 points)\ }
% solution macros
\newcommand{\ask}[1]{\par\textit{What is being asked.} #1\par}
\newcommand{\idea}[1]{\par\textit{The idea.} #1\par}
\newcommand{\pk}[1]{\par{\small\textit{[Packet: #1]}}\par}
\newcommand{\stp}[1]{\par\textbf{Step #1.}\ }
\newcommand{\ans}[1]{\[\boxed{\ #1\ }\]}
\newcommand{\trap}[1]{\par\textit{Watch out.} #1\par}
\newcommand{\sol}[1]{\par\bigskip\textbf{Part #1}\par}

\begin{document}

% ================================ EXAM ================================
\begin{center}
{\large\bfseries Macroeconomics I (ECO387C) Midterm Exam (Mock)}\\[1mm]
Fall 2026
\end{center}

Name: \rule{3in}{0.4pt}

\textbf{Instructions.} Closed book, no notes. Answer all four questions.
Show your derivations, since most of the credit is for the steps. Where a
number is asked for, leave it as a fraction or simple expression if you have
no calculator. Total 100 points, about 2 hours. Write your answers in the space below each part and continue on the back if needed.

% ---------------------------------------------------------------- Q1
\textbf{Question 1. A stochastic endowment economy} (25 points)

A continuum of identical households maximizes
$\E_0\sum_{t=0}^\infty\beta^t\frac{c_t^{1-\gamma}}{1-\gamma}$, $0<\beta<1$,
$\gamma>0$. Households trade a riskless one-period bond. Buying $a_{t+1}$
units at $t$ pays $R_ta_{t+1}$ units of goods at $t+1$, where $R_t$ is the
gross rate set at $t$. Bonds are in zero net supply. Every household
receives the same endowment $y_t\in\{y_h,y_l\}$ with $y_h>y_l>0$. The
endowment follows a Markov chain. From $y_h$ it stays at $y_h$ with
probability $p_h$, and from $y_l$ it stays at $y_l$ with probability $p_l$.

\needspace{6.0cm}
\qp{a}{2} Write the household's flow budget constraint.
\gap{3.5cm}
\needspace{7.0cm}
\qp{b}{3} What are the state variables of the household problem? Explain
why the current endowment is a state even though it is exogenous.
\gap{4.5cm}
\needspace{7.0cm}
\qp{c}{3} Write the Bellman equation, with the expectation written out as
a sum over next period's states.
\gap{4.5cm}
\needspace{10.0cm}
\qp{d}{5} Derive the Euler equation from the Bellman equation. Label the
first-order condition and the envelope condition.
\gap{7.5cm}
\needspace{8.5cm}
\qp{e}{4} State the market-clearing conditions. Use them to show that in
equilibrium $a_{t+1}=0$ and $c_t=y_t$, and give a formula for the
equilibrium gross rate $R(y)$.
\gap{6.0cm}
\needspace{8.5cm}
\qp{f}{4} Let $\gamma=2$, $y_h=2$, $y_l=1$. Find $R_h$ and $R_l$ as
functions of $\beta$, $p_h$, $p_l$. Evaluate them at $\beta=0.95$,
$p_h=0.8$, $p_l=0.6$ and compare with $1/\beta$.
\gap{6.0cm}
\needspace{8.5cm}
\qp{g}{4} With the same numbers, the log case ($\gamma=1$) gives
$R_h=0.877$ and $R_l=1.316$. Explain why the rates are further from
$1/\beta$ when $\gamma=2$. Then suppose the endowment is iid, $p_h=1-p_l$.
Show that $R_l/R_h$ no longer depends on the probabilities and find it.

\gap{6.0cm}
% ---------------------------------------------------------------- Q2
\newpage
\textbf{Question 2. An anticipated income rise} (25 points)

Time is $t=0,1,2,\dots$. There are two household types, $A$ and $B$, each
of mass one, with preferences $\sum_{t=0}^\infty\beta^t\ln c^e_t$ for
$e\in\{A,B\}$. Type $A$ receives $y^A_t=1$ every period. Type $B$ receives
$y^B_0=1$ and $y^B_t=2$ for all $t\ge1$. Everybody knows this at $t=0$.
Households trade one-period real bonds: $c^e_t+b^e_t=y^e_t+(1+r_{t-1})b^e_{t-1}$
with $b^e_{-1}=0$.

\needspace{8.5cm}
\qp{a}{4} State the household problem, including the no-Ponzi-game
condition, and define a sequential competitive equilibrium.
\gap{6.0cm}
\needspace{8.5cm}
\qp{b}{4} Derive the Euler equation and the intertemporal budget
constraint. In the intertemporal budget constraint, write
$q_t:=1/\prod_{\tau=0}^{t-1}(1+r_\tau)$.
\gap{6.0cm}
\needspace{8.5cm}
\qp{c}{4} Show that $1+r_t=\frac1\beta\frac{Y_{t+1}}{Y_t}$ with
$Y_t:=y^A_t+y^B_t$. Find $r_0$ and $r_t$ for $t\ge1$, and the prices $q_t$.
Explain why $r_0$ is high.
\gap{6.0cm}
\needspace{11.5cm}
\qp{d}{6} Solve for $c^A_t$ and $c^B_t$ for all $t$. Verify that the
goods market clears.
\gap{9.0cm}
\needspace{8.5cm}
\qp{e}{4} Find $b^A_0$ and $b^A_t$ for $t\ge1$. Who borrows at $t=0$,
and why? Check that $A$'s bond path satisfies the no-Ponzi-game
condition.
\gap{6.0cm}
\needspace{7.0cm}
\qp{f}{3} In the Arrow-Debreu version of this economy (textbook
Section 5.2.1), the price of date-$t$ goods in units of date-0 goods is
$p_t=\beta^tY_0/Y_t$. Show that $p_t=q_t$ and explain in one sentence why
the two market structures give the same allocation.

\gap{4.5cm}
% ---------------------------------------------------------------- Q3
\newpage
\textbf{Question 3. Growth by dynamic programming} (25 points)

A planner maximizes $\sum_{t=0}^\infty\beta^t\ln c_t$ subject to
$c_t+k_{t+1}=Ak_t^\alpha$, $k_{t+1}\ge0$, $k_0>0$ given, with
$0<\alpha<1$, $0<\beta<1$, $A>0$ (full depreciation).

\needspace{7.0cm}
\qp{a}{3} Write the problem in the form
$\max\sum\beta^tF(k_t,k_{t+1})$ with $k_{t+1}\in\Gamma(k_t)$. Give $F$ and
$\Gamma$, and write the Bellman equation.
\gap{4.5cm}
\needspace{10.0cm}
\qp{b}{5} Let $T$ be the Bellman operator. State Blackwell's sufficient
conditions and verify both for $T$. What does Banach's fixed-point theorem
then deliver? Why is some care needed with log utility?
\gap{7.5cm}
\needspace{10.0cm}
\qp{c}{5} Start value function iteration from $V_0\equiv0$. Find $V_1$
and its policy. Then find the policy that maximizes the right-hand side
when computing $V_2$, and interpret it as a decision in a finite-horizon
problem.
\gap{7.5cm}
\needspace{11.5cm}
\qp{d}{6} Derive the functional Euler equation using the first-order and
envelope conditions. Guess $g(k)=sAk^\alpha$ and solve for $s$.
\gap{9.0cm}
\needspace{7.0cm}
\qp{e}{3} Find the steady state $k^*$ and evaluate it at $\alpha=0.3$,
$\beta=0.95$, $A=1$.
\gap{4.5cm}
\needspace{7.0cm}
\qp{f}{3} Write the policy in log deviations
$\hat k_t:=\ln k_t-\ln k^*$. If $k_0=k^*/2$, what is $\hat k_1$? How many
periods does it take to close half of the gap?

\gap{4.5cm}
% ---------------------------------------------------------------- Q4
\newpage
\textbf{Question 4. Consumption} (25 points)

A household maximizes $\E_0\sum_{t=0}^\infty\beta^tu(c_t)$ subject to
$a_{t+1}=(1+r)(a_t+y_t-c_t)$ and a no-Ponzi-game condition, with
$\beta(1+r)=1$ and $r>0$. Income $y_t$ is stochastic.

\needspace{7.0cm}
\qp{a}{3} Let $u(c)=b_1c-\frac12b_2c^2$. Show that $c_t=\E_tc_{t+1}$.
\gap{4.5cm}
\needspace{10.0cm}
\qp{b}{5} Iterate the budget constraint forward and use part (a) to show
\[
c_t=\frac{r}{1+r}\Big[a_t+\sum_{j=0}^\infty\Big(\frac{1}{1+r}\Big)^j\E_ty_{t+j}\Big].
\]
\gap{7.5cm}
\needspace{11.5cm}
\qp{c}{6} Let $y_t-\bar y=\rho(y_{t-1}-\bar y)+\varepsilon_t$ with
$0\le\rho\le1$ and $\varepsilon_t$ white noise. Show that
$\Delta c_t=\frac{r}{1+r-\rho}\varepsilon_t$. Evaluate the coefficient at
$r=0.04$ for $\rho=0$, $0.9$, $1$ and interpret.
\gap{9.0cm}
\needspace{7.0cm}
\qp{d}{3} Hall (1978) regresses $\Delta c_{t+1}$ (or $\Delta\ln c_{t+1}$)
on variables known at $t$. What does the model predict for the coefficient
and why? What did Johnson, Parker and Souleles (2006) find?
\gap{4.5cm}
\needspace{10.0cm}
\qp{e}{5} Now add the borrowing constraint $a_{t+1}\ge0$ and let income be
iid with mean $\bar y=1$, so $\rho=0$. At $a_t=0$ the household draws
$y_t=0.5$, with $r=0.04$. Compute the consumption the unconstrained rule
from (b) prescribes, and show that it violates the constraint. Write the
Euler equation with the constraint's multiplier, and give consumption and
the marginal propensity to consume at this point.
\gap{7.5cm}
\needspace{7.0cm}
\qp{f}{3} Replace quadratic utility by CRRA. Explain, using a
second-order expansion or a picture, why the household now saves more
when income risk rises, even with no borrowing constraint.


\gap{4.5cm}

\end{document}
